% The local reduction ladder used by the compositional route.
%
% This section separates the additive target-count theorem, which does not
% need migration, from the full Kempe-reachability theorem, which does.

\section{The Meyniel ladder and the repaired square seam}
\label{sec:meyniel-ladder}

The reduction ladder removes a facial triangle, a digon, a two-edge cut, or
a facial square before the compositional argument reaches its pentagonal
frontier.  Two different invariants occur in the working manuscript:
nonnegative target counts, used by the Seed Lemma, and connectivity of the
Kempe reconfiguration graph, used by the stronger \(\mathrm{WC}^{*}\)
statement.  They agree on the first three rungs but diverge at a square.
The distinction is decisive: an additive identity transports a zero count
without any path between its two summands, whereas connectivity requires an
actual migration path.

\subsection{The three unbranched rungs}

For a facial triangle, properness forces its three internal edges to have
the three different colours.  At each corner the outer edge therefore has
the colour of the opposite triangle edge.  Contraction to one cubic vertex
is consequently canonical, and the three possible two-colour paths through
the contracted vertex expand to the three corresponding paths through the
triangle.  Both the colour law and the three path expansions are checked in
the rotation-system development.
\Sverified{Mettapedia.GraphTheory.FourColor.GoertzelV24Rotor.RotationSystem.rotor_outerColors_eq_oppositeTriangleEdges}
\Sverified{Mettapedia.GraphTheory.FourColor.GoertzelV24Rotor.RotationSystem.rotor_three_bicoloredThroughPaths}
The resulting finite portal profile is equal to the contracted-vertex
profile, including composition with an arbitrary exterior profile.
\Sverified{Mettapedia.GraphTheory.FourColor.GoertzelV24RotorProfile.RotationSystem.facialTriangle_composedRotorProfile_eq_contractedVertexProfile}

For a digon, its two parallel edges form one maximal two-colour component.
Switching that component exchanges the two local extensions, so the local
fibre is connected.  Suppression preserves the boundary connectivity
profile, with multiplicity two.
\Sverified{Mettapedia.GraphTheory.FourColor.GoertzelV24Digon.RotationSystem.digon_internalColorings_connected}
\Sverified{Mettapedia.GraphTheory.FourColor.GoertzelV24DigonProfile.digonProfileMultiplicityTwo}

Across a two-edge cut, the two crossing edges have equal colours.  This is
the Klein-group boundary-sum law on either side of the cut.  Capping the two
sides is exact in both directions: an ambient colouring restricts to cap
colourings, and cap colourings may be aligned by a global colour permutation
and glued.
\Sverified{Mettapedia.GraphTheory.FourColor.GoertzelV24TwoEdgeCut.RotationSystem.twoEdgeCut_colors_eq}
\Sverified{Mettapedia.GraphTheory.FourColor.GoertzelV24TwoEdgeCutCap.RotationSystem.TwoEdgeCutPairData.taitColorable_iff_caps}
The remaining work on these rungs is integration: package these local
equivalences as the exact target-count and reconfiguration transports
consumed by the common ladder interface.  No unproved geometric principle is
hidden in that packaging.

\subsection{The exact square fibre}

Let \(Q=wxyz\) be a facial square, and let
\((c_w,c_x,c_y,c_z)\) be the colours on its four outer edges in cyclic
order.  There are two planar reductions.  The first joins the outer edges at
\(w,x\) and at \(y,z\); the second joins those at \(x,y\) and at \(z,w\).
Write them as \(G_0,G_1\).

\begin{theorem}[square fibre table]
\label{thm:square-fibre-table}
\Sverified{Mettapedia.GraphTheory.FourColor.GoertzelV24Square.squareFiberTable}
For a nonzero boundary word, the square has two internal colourings when all
four boundary colours are equal, one internal colouring for either adjacent
pair pattern
\[
 (p,p,q,q)\quad\hbox{or}\quad(q,p,p,q),\qquad p\ne q,
\]
and no internal colouring otherwise.  In particular, the opposite-pair
pattern \((p,q,p,q)\) is impossible.
\end{theorem}

The two states in an all-equal fibre differ by switching the square's own
two-colour cycle.  Thus every nonempty local fibre is internally connected.
\Sverified{Mettapedia.GraphTheory.FourColor.GoertzelV24Square.squareFibers_internallyConnected}

For one exterior colouring profile \(\eta\), the two internal states over an
all-equal word have different through-square pairings: one belongs to the
\(G_0\) reduction and the other to \(G_1\).  An adjacent-pair word belongs to
exactly one reduction.  Keeping multiplicity therefore gives, for every
target interface profile \(\pi\),
\begin{equation}
  \operatorname{Count}_{G}(\eta,\pi)
  =\operatorname{Count}_{G_0}(\eta,\pi)
   +\operatorname{Count}_{G_1}(\eta,\pi).
  \label{eq:square-profile-additivity}
\end{equation}
This is proved for arbitrary multisets of exterior profiles, not merely for
one census graph.
\Sverified{Mettapedia.GraphTheory.FourColor.GoertzelV24SquareComposition.squareGluedAdditiveProfileCountIdentity}

\begin{lemma}[zero targets descend additively]
\label{lem:square-zero-descends}
\Sprose
Let \(S\) be any finite set of target profiles and put
\(N_H(S)=\sum_{\pi\in S}\operatorname{Count}_H(\pi)\).  If
Equation~\eqref{eq:square-profile-additivity} holds for every \(\pi\), then
\[
     N_G(S)=N_{G_0}(S)+N_{G_1}(S),
     \qquad
     N_G(S)=0\Longleftrightarrow
     N_{G_0}(S)=N_{G_1}(S)=0.
\]
\end{lemma}

\begin{proof}
Sum Equation~\eqref{eq:square-profile-additivity} over \(S\) and distribute
the finite sum.  All summands are natural numbers, so a sum of the two
totals is zero exactly when both totals are zero.
\end{proof}

Consequently the square rung for the Seed target is additive rather than
dynamical: if a square instance has no seed, neither reduction has a seed.
Here the counted state is pointed by a choice of majority pair: it is a pair
\((C,P)\) of a colouring and one of its two majority pairs.  Thus the two
pair sectors form a disjoint state space even when one colouring is a seed
for both pairs.  Positivity of its total count is exactly the existence of a
colouring and a majority pair witnessing the Seed conclusion; no
inclusion--exclusion assertion about overlapping sets of colourings is being
used.

\begin{theorem}[graph-level square count transport]
\label{prop:graph-square-count-transport}
\Sprose
For a facial square whose two planar reductions are valid tangles, the
profile-resolved count of pointed colourings satisfies
Equation~\eqref{eq:square-profile-additivity}.  Consequently the pointed Seed
count is additive and a least zero-Seed instance descends through both
smaller reductions.
\end{theorem}

\begin{proof}
Delete the four internal square edges but retain its four ordered outer
half-edges.  A global colouring is then uniquely a colouring of this common
exterior together with one member of the local square fibre over its boundary
word.  A colouring of either reduction is uniquely the same exterior
colouring together with the corresponding compatible planar pairing.
For a fixed exterior colouring, fixed majority pair, and fixed target
profile, the equality of these two multisets is exactly the checked local
profile identity after composition with that exterior profile.  Summing over
the finite multiset of exterior colourings proves the graph-level identity.

Now sum over the finite set of Seed profiles in the pointed state space and
apply Lemma~\ref{lem:square-zero-descends}.  If the total upstairs is zero,
both reduction totals are zero; since a facial-square reduction removes four
vertices, either valid reduction is a strictly smaller zero-Seed instance.
No migration statement occurs in the argument.
\end{proof}

\subsection{The universal \(A\ne\varnothing\) shortcut is false}

Addendum~XXVI asserts that, for any tangle in the ladder's scope and any two
disjoint cofacial edges \(E_1,E_2\) arising as the joined edges of a square
reduction, some Tait colouring has
\(\chi(E_1)=\chi(E_2)\).  Equivalently, the inverse square has a colouring in
its all-equal boundary fibre \(A\).  The following specimen refutes that
universal statement.

\begin{proposition}[forced-unequal square reduction]
\label{prop:a-nonempty-refuted}
\Srefuted{tools/fourcolor/v24_a_nonempty_counterexample.py}
There is a connected planar cubic tangle with two stubs on one face and two
vertex-disjoint cofacial edges \(E_1,E_2\) which are the joined edges of a
planar square reduction, such that the tangle has six Tait colourings and
\(\chi(E_1)\ne\chi(E_2)\) in all six.
\end{proposition}

\begin{proof}
Take internal vertices \(0,\ldots,5\), stub vertices \(6,7\), and edges
\[
\begin{split}
01,02,03,14,15,23,25,34,46,57.
\end{split}
\]
At the internal vertices use cyclic rotations
\[
\begin{array}{c|rrrrrr}
v&0&1&2&3&4&5\\ \hline
\rho_v&(1,2,3)&(0,4,5)&(3,0,5)&(4,0,2)&(1,3,6)&(2,1,7).
\end{array}
\]
The stub rotations are \((4)\) at \(6\) and \((5)\) at \(7\).  Directly
following the face permutation gives
\[
 (0,1,5,2),\quad (0,2,3),\quad (0,3,4,1),\quad
 (1,4,6,4,3,2,5,7,5).
\]
Hence \(V-E+F=8-10+4=2\), both stubs lie on the last face, and
\(E_1=01,E_2=34\) are disjoint edges of the face \((0,3,4,1)\).
Every non-boundary edge lies on one of the displayed ordinary cycles, so
there is no internal bridge.

In any Tait colouring, write
\(x=\chi(02),y=\chi(03),z=\chi(23)\).  The triangle \(023\) forces
\(x,y,z\) to be the three distinct colours.  Properness at vertex \(0\)
then forces \(\chi(01)=z\), while properness at vertex \(3\) forces
\(\chi(34)=x\).  Thus the marked colours are always different.
Conversely, once the ordered triple \((x,y,z)\) is chosen, properness forces
\[
 \chi(14)=y,\quad \chi(15)=x,\quad \chi(25)=y,
 \quad\chi(46)=\chi(57)=z,
\]
so each of the six permutations of the three colours gives one colouring.

Finally delete \(01,34\), insert the facial square
\(8,9,10,11\), and attach its vertices in cyclic order to
\(0,1,4,3\).  One planar reduction restores \(01,34\).  The other adds a
second copy of each of \(14,03\), which is a legitimate multigraph
reduction.  The displayed rotation extends to a spherical rotation with the
new square as a face; the complete rotations and face audit are recorded by
the cited executable fixture.  Thus the marked pair is exactly a square
reduction pair, as required.
\end{proof}

The failure is not caused by disconnected Kempe dynamics.  Exhaustion of
the inverse square gives six colourings in one Kempe class and no all-equal
boundary state.  The first reduction has six colourings in one class; the
parallel-edge reduction has none.  Its square identity is therefore
\(6=6+0\).  The specimen disproves the shortcut \(A\ne\varnothing\), but it
does not disprove the square conclusion.

The converse warning is equally important: the full universal SQ3
connectivity statement is also false, although a slightly larger specimen is
needed.

\begin{proposition}[universal SQ3 migration is false]
\label{prop:universal-sq3-refuted}
\Srefuted{tools/fourcolor/v24_sq3_migration_counterexample.py}
There is a connected planar cubic tangle with two stubs on one face and a
facial square such that each of its two planar reductions has a connected
Kempe reconfiguration graph, but the original tangle's reconfiguration graph
has two components.
\end{proposition}

\begin{proof}
Use the planar Frucht graph with edge set
\[
\begin{gathered}
01,06,07,12,17,23,28,34,39,45,49,56,5,10,6,10,\\
7,11,8,11,89,10,11.
\end{gathered}
\]
The disjoint edges \(06,34\) lie on the face
\((0,6,5,4,3,2,1)\).  Delete them, insert a square with vertices
\(12,13,14,15\), and attach those vertices in cyclic order to
\(0,6,4,3\).  Finally open the unrelated edge \(01\) into stub edges
\(0,16,1,17\).  The resulting tangle has sixteen cubic vertices and two
stubs.  The rotation cycles in the cited fixture give nine face orbits,
including the square \((12,13,14,15)\) and one face containing both stubs;
thus \(V-E+F=18-25+9=2\).  Deleting any non-boundary edge leaves the tangle
connected.

The finite audit assigns the three colours by backtracking, rejecting a
partial assignment as soon as two assigned edges at one cubic vertex agree.
It obtains exactly eighteen Tait colourings.  For every colouring and each
of the three colour pairs, it computes every maximal two-colour component,
switches that whole component, and identifies the resulting colouring in the
same eighteen-state table.  The resulting Kempe graph has components of
sizes six and twelve.  Every state in the six-state component has square word
of type \((p,p,q,q)\), and every state in the twelve-state component has the
other adjacent-pair type; no all-equal state exists.

The first planar reduction restores \(06,34\) and has six colourings in one
Kempe class.  The other joins \(64,30\) and has twelve colourings in one
Kempe class.  Its two lifted sides are exactly the two components just
described.  This proves the claimed failure.  All rotations, face walks,
colourings, and component switches are emitted and asserted by the cited
deterministic fixture.
\end{proof}

This second counterexample still does not refute the route's loaded
\(\mathrm{WC}^{*}\) target.  Its two stub colours agree in every one of the
eighteen states, so every state already has the extendable two-stub word.
The target reachability statement is true even though the stronger
single-class statement is false.  The repair must therefore retain the
actual target predicate instead of trying to recover universal
reconfiguration connectivity.

\subsection{The correct reachability obligation}

Let \(\mathcal K(G)\) be the graph whose vertices are the Tait colourings of
\(G\) and whose edges are legal maximal-component Kempe switches.  Let
\(L_0\) consist of the colourings whose square boundary is all-equal or of
the adjacent-pair type represented by \(G_0\), and define \(L_1\)
symmetrically.  The fibre table gives
\(\operatorname{Col}(G)=L_0\cup L_1\); their intersection is exactly the
all-equal set \(A\).

\begin{lemma}[connected-union criterion]
\label{lem:square-connected-union}
\Sprose
Suppose each nonempty induced reconfiguration subgraph on \(L_0,L_1\) is
connected.  Then \(\mathcal K(G)\) is connected precisely when one of the
sets is empty or some Kempe path has one endpoint in \(L_0\) and the other
in \(L_1\).
\end{lemma}

\begin{proof}
Necessity is immediate.  For sufficiency, join any colouring first inside
its own connected side to the appropriate endpoint of the cross-side path,
traverse that path if the target lies on the other side, and then join to
the target inside that side.  If one side is empty, the other is the whole
colouring set.
\end{proof}

Full connectivity is stronger than the route needs.  For a target set
\(W\) of stub words, put
\(W_s=\{C\in L_s:\operatorname{word}(C)\in W\}\).

\begin{lemma}[target-aware square criterion]
\label{lem:target-aware-square}
\Sprose
Assume that, whenever \(W_s\ne\varnothing\), every state of \(L_s\) can
reach a state of \(W_s\) while staying in \(L_s\).  If the global target set
is nonempty, then \(\mathrm{WC}^{*}(G,W)\) holds if and only if every state
in a target-free side \(L_s\) can reach some side \(L_t\) with
\(W_t\ne\varnothing\).
\end{lemma}

\begin{proof}
For necessity, a path from a target-free side to a global target has a final
state in some target-bearing side.  Conversely, a state already in a
target-bearing side reaches \(W\) by the assumed sidewise statement.  A
state in a target-free side first takes the supplied path to a
target-bearing side and then uses its sidewise path to \(W\).
\end{proof}

The target cannot be left abstract in the remaining square theorem.
Even colour-permutation invariance, planarity, bridgelessness, and one-face
placement of the stubs do not make generic \(\mathrm{WC}^{*}\) transport true.

\begin{proposition}[generic \(\mathrm{WC}^{*}\) square transport is false]
\label{prop:generic-wcstar-square-refuted}
\Srefuted{tools/fourcolor/v24_square_wcstar_counterexample.py}
There are a bridgeless planar cubic tangle \(G\) with four stubs on one face,
a facial square with two bridgeless planar reductions \(G_0,G_1\), and a
colour-permutation-invariant set \(W\) of stub words such that
\[
  \mathrm{WC}^{*}(G_0,W)\quad\hbox{and}\quad
  \mathrm{WC}^{*}(G_1,W)
  \qquad\hbox{but}\qquad
  \neg\mathrm{WC}^{*}(G,W).
\]
\end{proposition}

\begin{proof}
Start with the same planar Frucht rotation used in
Proposition~\ref{prop:universal-sq3-refuted}.  Replace its cofacial edges
\(06,34\) by the square \(12\,13\,14\,15\), attached in cyclic order to
\(0,6,4,3\).  Instead of opening \(01\), open the two exterior edges
\(5\,10\) and \(8\,11\).  Denote the resulting stub edges, in that order, by
\[
  5\,16,\quad 10\,17,\quad 8\,18,\quad 11\,19.
\]
Let \(W\) consist of the words
\[
             (a,a,b,b),\qquad a,b\in\{0,1,2\},\quad a\ne b.
\]
This set is invariant under every permutation of the three colours.

The rotation is obtained from the closed Frucht rotation by replacing the
neighbour \(10\) at \(5\) by \(16\), the neighbour \(5\) at \(10\) by
\(17\), the neighbour \(11\) at \(8\) by \(18\), and the neighbour \(8\)
at \(11\) by \(19\).  Its face permutation has eight orbits, including
\((12,13,14,15)\) and a single orbit containing \(16,17,18,19\).
Thus \(V-E+F=20-26+8=2\), the square is facial, and the four stubs lie on
one face.  The two reductions respectively restore \(06,34\) and insert
\(64,30\).  Each has \(16\) vertices, \(20\) edges, six face orbits, and one
face containing all four stubs.  Deleting each internal edge in turn leaves
all vertices connected in all three maps, so all three tangles are
bridgeless.  The complete rotations and face walks are asserted in the
cited fixture.

It remains only a finite colouring calculation.  Assign colours to indexed
edges by backtracking, rejecting an assignment as soon as two assigned edges
at a cubic vertex agree.  For each complete colouring and each of the three
colour pairs, find every maximal two-colour component and switch the whole
component.  Looking the result up in the complete colouring table gives the
exact Kempe graph.  This calculation yields
\[
\begin{array}{c|c|c|c}
 &|\operatorname{Col}|&\text{Kempe-class sizes}&\text{\(W\)-states per class}\\
 \hline
 G   &48&18,\ 30&6,\ 0\\
 G_0 &30&30&0\\
 G_1 &18&18&6.
\end{array}
\]
Every search is exhaustive; there is no random or heuristic step.

By the source definition, \(\mathrm{WC}^{*}(G_0,W)\) is true because no
colouring of \(G_0\) has a word in \(W\), so its proviso is false.
The graph \(\mathcal K(G_1)\) is connected and contains six target states,
so \(\mathrm{WC}^{*}(G_1,W)\) is true.  Upstairs the target is nonempty but
the thirty-state Kempe class contains none of it.  Hence
\(\mathrm{WC}^{*}(G,W)\) is false.
\end{proof}

This proposition does not refute the route's surviving square claim.  Its
target \(W\) is not the pentagonal cap's extendable-word set, and the
specimen is not asserted to be a least obstruction.  It proves that both of
those qualifiers carry mathematical content: no theorem quantified over an
arbitrary colour-invariant target can replace them.

The local square transparency theorem supplies the expected sidewise
lifting mechanism.  A colouring of \(G_s\) expands to \(L_s\); its local
fibre is connected by the square-cycle switch; and a Kempe component through
a joined edge expands to the corresponding boundary-to-boundary path through
the square.  Establishing this simultaneously for every ambient maximal
component is an exact graph-level path-lifting theorem, distinct from the
finite local fibre computation.

\begin{theorem}[square-side path lifting]
\label{prop:square-side-path-lifting}
\Sprose
For each reduction side \(s\), projection
\(L_s\to\operatorname{Col}(G_s)\) is surjective, has connected fibres, and
lifts every legal Kempe edge of \(G_s\) to a path in the induced
reconfiguration graph on \(L_s\).  Hence connectedness of
\(\mathcal K(G_s)\) implies connectedness of \(L_s\).
\end{theorem}

\begin{proof}
Fix a reduction colouring.  Give each of the two outer half-edges replacing
a joined edge the colour of that joined edge.  The resulting square boundary
word is compatible with side \(s\).  The fibre table therefore supplies a
lift.  An adjacent-pair word has one lift; an all-equal word has two, and the
two are joined by the square-cycle Kempe move.  Projection is consequently
surjective with connected fibres.

Consider one downstairs Kempe move, switching colours \(a,b\) on a maximal
component \(K\).  For the present boundary word, the checked local profile
identity says that the multiset of \((a,b)\)-portal connectivity relations of
the square extensions is the multiset of the compatible reduction pairings.
If necessary, perform the square-cycle fibre move first; one of the lifts now
has exactly the portal relation of side \(s\) for \((a,b)\).

Glue the common exterior to the local tracked-colour graph.  Components of
the glued graph which meet the exterior are obtained by taking exterior
components and identifying their incident portals according to the local
portal relation.  Equal portal relations therefore give a bijection between
the downstairs maximal components and the upstairs maximal components which
meet the exterior, preserving every exterior edge.  In particular \(K\)
lifts to one maximal component \(\widetilde K\).  Switch \(a,b\) on
\(\widetilde K\).  At the four square vertices this is an ordinary Kempe
switch on a maximal component, so properness is preserved; outside the square
its support is exactly the support of the downstairs switch.  The resulting
colouring therefore projects to the prescribed successor of the reduction
colouring and remains compatible with side \(s\).

Thus one edge of \(\mathcal K(G_s)\) lifts to at most two upstairs edges: an
optional internal square-cycle move followed by the lifted component move.
Lift a path downstairs one edge at a time, and use connectedness of the final
fibre to reach any chosen lift of its endpoint.  This proves the last
assertion.
\end{proof}

The source target supplies one further fact which an arbitrary set of words
does not: a target word is exactly the restriction of a colouring after the
fixed completion cap is restored.  Minimality can therefore prove that each
reduction side contains a target state.

\begin{theorem}[route-specific square transport from cap nonemptiness]
\label{prop:target-square-migration}
\Sprose
Let \(G\) be a tangle with a facial square disjoint from its ordered boundary,
let \(G_0,G_1\) be its two valid planar reductions, and let \(C\) be the
fixed completion cap at that boundary.  Let \(W_C\) be the set of boundary
words which extend across \(C\).  Assume
\begin{enumerate}
\item each capped reduction \(G_s\cup C\) is Tait-colourable;
\item \(\mathrm{WC}^{*}(G_s,W_C)\) holds for \(s=0,1\); and
\item square-side path lifting
  (Theorem~\ref{prop:square-side-path-lifting}) holds.
\end{enumerate}
Then \(\mathrm{WC}^{*}(G,W_C)\) holds.

In particular, a facial square cannot occur in a least obstruction of the
source's strong induction whenever attaching the same cap to either
reduction gives a strictly smaller closed instance: closed-instance
minimality supplies the first assumption and the tangle induction supplies
the second.
\end{theorem}

\begin{proof}
Restrict a Tait colouring of \(G_s\cup C\) to \(G_s\).  Its boundary word
extends across \(C\) by the colouring just removed, so
\[
 W_s=\{D\in\operatorname{Col}(G_s):
                  \operatorname{word}(D)\in W_C\}
\]
is nonempty for both \(s=0,1\).

Now take any colouring \(D\) of \(G\).  The square fibre table puts \(D\) in
at least one lifted side \(L_s\).  Project it to a colouring
\(\overline D\) of \(G_s\).  Since \(W_s\ne\varnothing\), the assumption
\(\mathrm{WC}^{*}(G_s,W_C)\) supplies a Kempe path from
\(\overline D\) to some \(\overline E\in W_s\).  Lift that path one edge at
a time by Theorem~\ref{prop:square-side-path-lifting}; its proof permits an
optional square-cycle move inside a fibre and then the corresponding maximal
component move.  The lifted endpoint is a lift \(E\) of \(\overline E\).
Projection changes only the square and preserves every ordered boundary
edge, hence
\(\operatorname{word}(E)=\operatorname{word}(\overline E)\in W_C\).
Thus every colouring of \(G\) reaches the target.

The target of \(G\) is nonempty as well: choose either
\(\overline E\in W_s\) and use the surjectivity part of square-side lifting
to lift it.  Therefore the proviso in the definition of
\(\mathrm{WC}^{*}\) is respected, and the claimed statement follows.
\end{proof}

This proof identifies precisely why the generic counterexample is harmless.
There \(G_0\) has no target colouring, so its capped-colourability premise
fails; the vacuous truth of \(\mathrm{WC}^{*}(G_0,W)\) cannot create a
target.  In the source induction, target nonemptiness comes from the smaller
capped instance, not from an all-equal square state.  Hence neither
universal \(A\ne\varnothing\), universal SQ3 migration, nor arbitrary-target
transport may be restored.

Thus the square ledger is now exact:
\[
\begin{array}{rcl}
\text{local fibre table and fibre connectivity}&:&\text{machine-checked},\\
\text{profile-resolved additive identity}&:&\text{machine-checked},\\
\text{Seed zero-count descent}&:&\text{ordinary proof; Lean formalization pending},\\
\text{universal }A\ne\varnothing&:&\text{refuted},\\
\text{universal SQ3 migration}&:&\text{refuted},\\
\text{generic target-aware }\mathrm{WC}^{*}\text{ transport}&:&\text{refuted},\\
\text{route-specific }\mathrm{WC}^{*}\text{ square transport}&:&
   \text{ordinary proof; Lean formalization pending}.
\end{array}
\]
